\documentclass{article}

% Use the official ICLR 2027 style when it is present. The fallback keeps this
% source readable and locally compilable while the official style bundle is
% unavailable.
\IfFileExists{iclr2027_conference.sty}{
  \usepackage{iclr2027_conference,times}
}{
  \usepackage[margin=1in]{geometry}
  \usepackage{times}
}

\usepackage[T1]{fontenc}
\usepackage[utf8]{inputenc}
\usepackage{microtype}
\usepackage{graphicx}
\usepackage{booktabs}
\usepackage{amsmath,amssymb,amsfonts}
\usepackage{mathtools}
\usepackage{enumitem}
\usepackage{xcolor}
\usepackage{url}
\usepackage{natbib}
\usepackage{subcaption}
\usepackage{hyperref}
\usepackage[nameinlink,noabbrev]{cleveref}

\newcommand{\resultpanel}[1]{%
  \IfFileExists{#1}{%
    \includegraphics[width=\linewidth]{#1}%
  }{%
    \fbox{\parbox[c][1.05in][c]{0.92\linewidth}{%
      \centering\small Figure asset not present in this source bundle.}}%
  }%
}

\title{Routing-Conditioned Optimization Shapes Sparse Supports in
Mixture-of-Experts Models}
\author{Anonymous Authors\\Paper under double-blind review}
\date{}

\begin{document}

\maketitle

\begin{abstract}
Token routing in sparse Mixture-of-Experts (MoE) models determines which
gradients each expert receives, so some routing dependence of expert parameters
is expected. We ask whether this dependence persists as structured,
optimization-relevant support within expert weights. We use expert-local
magnitude masks as a probe in balanced decoder-only MoEs trained on a
50k-story TinyStories subset, a WikiText-103 subset exported from the first
10k source rows, and a
character-balanced mixture, and test the masks with matched controls,
rewinding, and interventions on routed-token history. First, conventionally
trained experts contain non-random, expert-specific sparse structure: in the
reference WikiText setting, 50\%-sparse learned masks increase validation loss
by 3.1\%, versus more than 140\% for matched random masks. Second, 50\%-sparse
supports retrain successfully from initialization in the evaluated settings,
whereas 80\%-sparse supports generally require early rewinding, showing that
topology at aggressive sparsity is not sufficient. Third, exact primary-ID
replay, random routing, and utilization-preserving token shuffling show that routed-token
history---including token identity and gate/capacity effects---affects support
formation beyond the logged primary-selection counts.
Under a new initialization, the recorded source primary-expert trajectory is
reproduced exactly but does not act as a portable sparse blueprint. These
findings support a bounded conclusion:
routing history makes a causal contribution within the tested regimes, while
expert-local sparse subnetworks emerge through the interaction of routing,
initialization, and optimization.
\end{abstract}

\section{Introduction}
\label{sec:introduction}

Sparse Mixture-of-Experts (MoE) models expand parameter capacity while routing
each token through only a small subset of expert modules
\citep{shazeer2017outrageouslylargeneuralnetworks,
lepikhin2020gshardscalinggiantmodels,
fedus2022switchtransformersscalingtrillion}. Conditional computation also
changes training: a routing decision determines which expert receives the
gradient associated with a token. Experts with identical architectures can
therefore experience different ordered sequences of optimization signals as
the router and experts co-adapt.

Because routing determines which gradients each expert receives, some
dependence of expert parameters on routing is expected by construction. Our
question is not whether routing changes expert weights, but whether it produces
persistent, structured, and optimization-relevant sparse supports that remain
detectable under controlled interventions and rewind-and-retrain tests.

Prior work has characterized routing stability, load balancing, expert
function, and redundancy, while MoE pruning shows that trained experts contain
removable parameters
\citep{dai-etal-2022-stablemoe,
guo2026advancingexpertspecializationbetter,
xie2024moeprunerpruningmixtureofexpertslarge}. Neither structured routing nor
successful final-checkpoint pruning, however, establishes that routed
optimization leaves expert-specific sparse structure with continued
optimization significance. Post-training compression can preserve a trained
function without identifying a subnetwork that remains trainable after its
weights are rewound.

We use expert-local sparse supports as a probe of this unresolved mechanism.
Within each expert, we extract magnitude masks and compare them with matched
random masks, masks transferred from other experts, and randomly reinitialized
surviving weights. We then rewind the masked models to initialization or early
checkpoints and retrain them, following the distinction between classical
initialization-rewind tickets and early-rewind behavior
\citep{frankle2019lotterytickethypothesisfinding,pmlr-v119-frankle20a}.
Finally, we intervene on routing through replay, random assignment,
selection-count-preserving token shuffling, and cross-initialization replay.
These tests contrast realized routed-token histories at matched logged
primary-selection counts and test whether routing alone determines the
resulting support. Because the
archive does not store which tokens survive the capacity filter or their
gate-weighted gradient mass, the shuffle is not an exact match of realized
expert updates.

The evidence supports three claims. First, conventionally trained MoE experts
contain non-random, expert-specific sparse subnetworks. Second, these supports
retain optimization significance: moderate-sparsity supports satisfy the
initialization-rewind criterion in the evaluated settings, whereas aggressive
sparsity generally requires early optimization. Third, routing history makes a
causal contribution to support formation beyond utilization alone, but the same
routing trajectory need not produce the same sparse outcome under another
initialization. We therefore interpret sparse-support formation as an
interaction among routing history, initialization, and optimization, not as a
consequence of routing alone. Our claims are restricted to the small, balanced
decoder-only MoEs and interventions studied here.

\paragraph{Contributions.}
\begin{itemize}[leftmargin=*]
  \item \textbf{Structural.} We show that conventionally trained MoE experts
  contain non-random, expert-specific sparse supports that preserve global and
  routed-token performance substantially better than matched random,
  transferred, and reinitialized controls.

  \item \textbf{Optimization.} We distinguish final-checkpoint compression
  from lottery-ticket behavior through rewind-and-retrain tests. At 50\%
  sparsity, learned supports satisfy the initialization-rewind criterion in the
  evaluated settings; at 80\%, they generally require an early rewind and
  continue to outperform matched sparse controls.

  \item \textbf{Routing conditioning.} Through controlled routing
  interventions, we show that routed-token history contributes to
  sparse-support formation beyond logged primary-selection counts.
  Cross-initialization
  replay further rejects the hypothesis that routing independently determines
  the support, supporting an interaction among routing, initialization, and
  optimization.
\end{itemize}

\section{Related Work}
\label{sec:related_work}

\paragraph{MoE routing and expert specialization.}
Sparse MoEs use learned routing to allocate tokens to a small subset of experts
\citep{shazeer2017outrageouslylargeneuralnetworks,
lepikhin2020gshardscalinggiantmodels,
fedus2022switchtransformersscalingtrillion}. Work on routing fluctuation and
load balancing shows that assignments alter experts' training exposure
\citep{dai-etal-2022-stablemoe,
guo2026advancingexpertspecializationbetter}. Analyses of expert function and
representation geometry further caution that structured routing need not imply
a human-interpretable semantic role
\citep{lo2025closer,wang2026myth,tang-etal-2026-specialization}. We therefore
study expert-local parameter differentiation rather than claiming semantic
specialization.

\paragraph{MoE pruning and learned modular structure.}
MoE compression removes complete experts or prunes weights within surviving
experts using parameter, activation, or routing-derived importance
\citep{chen2022taskspecific,lu2024notall,
xie2024moeprunerpruningmixtureofexpertslarge,
lee2025stunstructuredthenunstructuredpruningscalable}. Related approaches
construct conditional sparse subnetworks from predefined data partitions or
learn modular topology from gradient conflict
\citep{stefanski2026routing,gan2025cdsp}. These results establish redundancy
and the usefulness of sparse modularity, but they do not test whether ordinary
MoE routing histories shape expert-local supports or whether those supports
remain trainable after rewinding.

\paragraph{Lottery tickets and the remaining gap.}
The Lottery Ticket Hypothesis identifies sparse supports that train from their
original initialization, while later work shows that difficult settings may
require rewinding to an early checkpoint
\citep{frankle2019lotterytickethypothesisfinding,pmlr-v119-frankle20a}.
Mask identity can contain optimization-relevant information beyond sparsity
alone \citep{paul2023unmasking}. We retain these criteria: routing history is a
candidate mechanism of ticket formation, not an added component of the formal
ticket definition. Prior work studies routing behavior, expert function, or
MoE compression; we test whether training-time routing shapes expert-local
sparse supports that retain optimization significance under rewinding.

\section{Routing-Conditioned Sparse-Subnetwork Formation}
\label{sec:formation}

\subsection{Problem and Hypotheses}
\label{sec:problem}

Let $\Theta_0$ be the initialization of the complete MoE and let
$a_{s,b,t,e}^{(\ell)}\in\{0,1\}$ indicate that expert $e$ is the recorded
primary pre-capacity selection for token $(b,t)$ in layer $\ell$ at training
step $s$. The archived routing trajectory, ordered selected-token history, and
aggregate primary-selection count are
\begin{equation}
\mathcal{A}
=
\left(a_{s,b,t,e}^{(\ell)}\right)_{\ell,s,b,t,e},
\qquad
\mathcal{H}_{\ell,e}
=
\operatorname{seq}
\left\{
(s,b,t,x_{s,b,t}):a_{s,b,t,e}^{(\ell)}=1
\right\},
\qquad
U_{\ell,e}
=
\sum_{s,b,t}a_{s,b,t,e}^{(\ell)}.
\label{eq:routing_objects}
\end{equation}
Two conditions can therefore match $U_{\ell,e}$ while assigning different
token identities and ordered histories. Capacity acceptance and raw gate values
mediate which selected tokens actually contribute to an expert update and are
not part of the archived trajectory.

After dense training, expert-local magnitude pruning extracts
\begin{equation}
m_{\ell,e}^{(\rho)}
=
\mathcal{P}_{\rho}\!\left(\theta_{\ell,e,T}\right),
\label{eq:expert_mask}
\end{equation}
where $\mathcal{P}_{\rho}$ retains the largest $1-\rho$ fraction of prunable
weights independently within each expert. Let $M^{(\rho)}$ be the corresponding
full-model mask, equal to one outside the prunable expert weight matrices. At
rewind point $t_0$, we retrain
\begin{equation}
\Theta'_{T;t_0}
=
\operatorname{Retrain}
\left(
M^{(\rho)}\odot\Theta_{t_0},
M^{(\rho)},
\mathcal{A}^{\mathrm{ret}}
\right),
\qquad
\Delta_{\mathrm{ret}}
=
\frac{
\mathcal{L}_{\mathrm{val}}(\Theta'_{T;t_0})
-
\mathcal{L}_{\mathrm{dense}}
}{
\mathcal{L}_{\mathrm{dense}}
}.
\label{eq:rewind_criterion}
\end{equation}
We describe a support as near dense when
$\Delta_{\mathrm{ret}}\leq\varepsilon=0.05$ and report its continuous gap and
gaps to matched random-mask, random-reinitialization, and randomized-routing
controls. The case $t_0=0$ tests the classical initialization-rewind criterion;
$t_0>0$ tests early-rewind behavior.

The experiments distinguish two hypotheses:
\begin{align}
\mathcal{H}_{\mathrm{route}}:\quad
m_{\ell,e}^{(\rho)}
&=
g_{\ell,e}(\mathcal{A}),
\label{eq:routing_only}\\
\mathcal{H}_{\mathrm{int}}:\quad
m_{\ell,e}^{(\rho)}
&=
g_{\ell,e}(\mathcal{A},\Theta_0,\mathcal{O}),
\label{eq:interaction}
\end{align}
where $\mathcal{O}$ collects the ordered data, optimizer, schedule, objective,
model interactions, and remaining optimization stochasticity. The routing-only
hypothesis predicts that an imposed trajectory determines support independently
of initialization. The interaction hypothesis predicts that changing routing
at fixed initialization changes support formation, while replaying the same
trajectory under another initialization need not reproduce the same outcome.

\subsection{Routing-Conditioned Optimization}
\label{sec:routing_framework}

Our working mechanism is narrower than a general theory of expert
specialization. Routing determines which token losses contribute to an
expert's updates; the resulting gradient trajectory interacts with current
parameters and optimization to shape the final weights and sparse support:
\begin{equation}
(\mathcal{A},\Theta_0,\mathcal{O})
\longrightarrow
\underbrace{
\left(\nabla_{\theta_{\ell,e}}\mathcal{L}_s\right)_{s=1}^{T}
}_{\substack{\text{capacity- and gate-mediated}\\\text{gradient trajectory}}}
\longrightarrow
\underbrace{
\left(\Theta_s\right)_{s=1}^{T}
}_{\text{parameter trajectory}}
\longrightarrow
m^{(\rho)}
\overset{\text{rewind/retrain}}{\longrightarrow}
\Delta_{\mathrm{ret}}.
\label{eq:mechanistic_chain}
\end{equation}
Primary-selection counts alone cannot specify this trajectory because gradient
values depend on token identity, order, gate values, capacity acceptance, and
the contemporaneous model state. Conversely,
routing alone need not specify it because gradients also depend on $\Theta_0$
and subsequent optimization. Sparse masks are therefore a probe of persistent
routing-conditioned structure, and rewindability tests whether that structure
retains optimization significance.

\subsection{Experimental Design}
\label{sec:experimental_design}

We study decoder-only Transformer MoEs trained by next-token prediction on a
50k-story TinyStories subset \citep{eldan2023tinystories}, a WikiText-103 subset
exported from the first 10k source rows \citep{merity2017pointer}, and a character-balanced mixture
of the two. The reference
model has four layers, four attention heads, width 256, eight experts per layer,
expert hidden width 1024, and top-1 routing with capacity factor 1.25. Runs use
sequences of 128 tokens, batches of 128 sequences, 2,500 updates, a
load-balancing coefficient of 0.1, and seeds 7, 17, and 29. The robustness suite
varies experts, routing width, depth, dataset, tokenizer, balancing strength,
and budget (\cref{fig:design}).

\begin{figure*}[t]
\centering
\begin{minipage}[t]{0.56\textwidth}
\centering
\small
\textbf{Mechanistic chain}\\[3pt]
\fbox{\parbox{0.94\linewidth}{\centering
$\mathcal{A}$ (routed-token history)
$\longrightarrow$ gradient trajectory
$\longrightarrow$ parameter trajectory
$\longrightarrow m^{(\rho)}
\longrightarrow$ rewindability}}
\vspace{5pt}

\begin{tabular}{@{}p{0.28\linewidth}p{0.65\linewidth}@{}}
\toprule
Test & Question isolated \\
\midrule
Sparse controls & Is the learned support non-random and expert-specific? \\
Rewind/retrain & Does the support retain optimization significance? \\
Primary-ID replay & Is the intervention faithful at fixed initialization? \\
Shuffled usage & Does realized routing matter beyond assignment count? \\
Cross-init replay & Does routing determine support independently of initialization? \\
\bottomrule
\end{tabular}
\end{minipage}
\hfill
\begin{minipage}[t]{0.40\textwidth}
\centering
\small
\textbf{Reference setting and extensions}\\[3pt]
\begin{tabular}{@{}p{0.28\linewidth}p{0.67\linewidth}@{}}
\toprule
Reference & 4 layers, 4 heads, $d=256$ \\
Experts & $E=8$, top-1, hidden 1024 \\
Data & TinyStories subset; WikiText-103 subset; mixture \\
Sequence / batch & 128 / 128 \\
Training & 2,500 steps; $\lambda_{\rm bal}=0.1$ \\
Seeds & 7, 17, 29 \\
Sparsity & 50\%, 80\% principal \\
Rewinds & 0\%, 1\%, 5\%, 10\% \\
\midrule
Extensions & $E\in\{4,8,16\}$; top-1/2 \\
 & 4/8 layers; byte-ngram tokenizer \\
 & representative 10,000-step suite \\
\bottomrule
\end{tabular}
\end{minipage}
\caption{\textbf{Question, mechanism, and design.} Sparse masks are used as a
probe of persistent structure rather than as a definition of expert function.
The intervention suite separates non-random support and rewindability, and
contrasts realized routing histories at matched utilization and initialization.
Full implementation and metric
definitions appear in \cref{app:setup,app:metrics}.}
\label{fig:design}
\end{figure*}

We jointly rank the two expert weight matrices within every expert and prune
the resulting expert-local vector; routers, attention, embeddings,
normalization parameters, output parameters, and expert biases remain dense.
Direct evaluation compares learned magnitude masks with matched random masks, a
limited transfer control that replaces expert 1's mask with expert 0's mask in
each layer, and learned masks applied to randomly reinitialized weights. The
primary metric is global validation loss relative to the dense model from the
same routing condition. Recorded expert-local losses use each evaluated
model's own selected routes; they are secondary diagnostics, not
fixed-assignment comparisons.

Rewind experiments evaluate 50\% and 80\% sparsity from initialization and
early checkpoints, keep pruned coordinates at zero, and run every condition for
a fresh full training budget from the same deterministically restarted data
order. The routing suite comprises same-initialization primary-ID replay,
random-every-step routing, selection-count-preserving shuffled routing, a
frozen random router, expert-trajectory permutations, and cross-initialization
replay. Except when initialization is the explicit intervention, paired
conditions share architecture, initialization, ordered token stream,
optimization settings, and training budget. Results are paired by training seed
and summarized as mean $\pm$ one standard deviation.
\Cref{app:setup,app:metrics,app:interventions} provide the complete protocol
available in the experimental record.

\section{Results}
\label{sec:results}

Unless otherwise noted, results are averaged over seeds 7, 17, and 29, with
variability reported as one standard deviation. Pruning comparisons use the
dense model trained under the same routing condition as their reference.

\subsection{Experts Contain Non-Random Sparse Supports}
\label{sec:results_sparse_structure}

Expert-local magnitude masks preserve substantially more of the trained model's
performance than matched sparse controls. On the WikiText-103 subset, pruning 50\% of each
expert's weights increased validation loss by 3.1\%, compared with increases
exceeding 140\% for random masks and 540\% when the surviving weights were
randomly reinitialized (\cref{fig:direct_pruning}). The corresponding
learned-mask increase on TinyStories was 4.6\%. Across the reported datasets and
architecture variants, 50\% magnitude pruning increased validation loss by at
most approximately 8\%. Learned masks also outperformed masks transferred
between experts in the limited expert-0-to-expert-1 control at the principal
50\% and 80\% sparsities in the original suites.

The separation from controls persists at 80\% sparsity, but direct pruning at
this level no longer retains near-dense performance. Thus, final checkpoints
contain non-random, expert-specific sparse organization, but direct-pruning
performance alone establishes neither rewindability nor semantic
specialization. The learned-routing architecture-grid aggregates also finished
with zero dead experts and normalized selected-route usage entropy between
0.9846 and 0.9952. These diagnostics rule
out gross utilization collapse, not functional similarity among experts.

\begin{figure*}[t]
\centering
\begin{subfigure}[t]{0.48\textwidth}
  \centering
  \resultpanel{figures/paper/fig2a_direct_pruning.pdf}
  \caption{Paired direct-pruning loss ratios on the WikiText-103 subset.}
  \label{fig:direct_pruning}
\end{subfigure}
\hfill
\begin{subfigure}[t]{0.48\textwidth}
  \centering
  \resultpanel{figures/paper/fig2b_rewind_80.pdf}
  \caption{Paired 80\%-sparse rewind gaps on the WikiText-103 subset.}
  \label{fig:rewind_pruning}
\end{subfigure}
\caption{\textbf{Sparse structure and optimization significance.}
(a) Per-seed loss is divided by the paired dense loss; the base-2 log scale
keeps the magnitude/control separation visible. Faint points are individual
seeds. Learned magnitude masks retain substantially more final-checkpoint
performance than matched random, limited 0-to-1 transfer, and reinitialization
controls. (b) At 80\% sparsity, rewinding to the 5--10\% checkpoints reduces
the learned-mask gap to the paired dense model and remains better than the
sparse controls. The dotted
line marks the 5\% ticket tolerance. Large points and curves are three-seed
means; error bars denote $\pm1$ sample standard deviation.}
\label{fig:structure_rewind}
\end{figure*}

\subsection{Sparse Supports Retain Optimization Significance Under Rewinding}
\label{sec:results_rewinding}

At 50\% sparsity, learned expert-local supports retrained successfully from
initialization on both the WikiText-103 and TinyStories subsets and outperformed matched
random-mask, random-reinitialization, and randomized-routing controls. These
experiments satisfy the initialization-rewind lottery-ticket criterion in the
evaluated settings. We reserve the term \emph{lottery ticket} for this
rewind-and-retrain evidence rather than for direct final-checkpoint pruning.

At 80\% sparsity, initialization rewind was generally insufficient, while early
rewinding substantially improved retraining (\cref{fig:rewind_pruning}).
In the three-seed TinyStories aggregate, a learned mask rewound to 10\% of
training reached validation loss $1.0413\pm0.0138$, 2.03\% above the dense loss
of $1.0206\pm0.0096$. At the same rewind point, the random-mask,
random-reinitialization, and randomized-routing controls reached
$1.1606\pm0.0035$, $1.1534\pm0.0115$, and $1.4579\pm0.0209$, respectively;
direct evaluation of the final mask without retraining yielded
$3.5043\pm0.3249$. These
results distinguish the learned support from a post-training compression
artifact: moderate-sparsity supports can remain trainable from initialization,
whereas the reported high-sparsity behavior requires a compatible early
parameter state and optimization trajectory.

\subsection{Routed-Token History Shapes Sparse-Support Formation}
\label{sec:results_routing}

We intervene on routing while matching initialization, data order,
architecture, and optimization settings. Same-initialization replay is the
positive control: the recorded result reports route agreement and mask
similarity of 1.0/1.0 and dense performance within experimental variability.

Disrupting the token-to-expert trajectory weakens the functional advantage of
learned masks. Random-every-step routing removes a coherent assignment history,
while shuffled-usage routing preserves the logged number of primary selections
for each expert at every layer and training step but permutes the assigned token
identities (\cref{fig:routing_pruning}). The learned-mask advantage is
substantially reduced under both interventions. Because pruning is normalized
to the condition-specific dense baseline, this is not solely a consequence of
different dense-model quality. In particular, shuffled usage rejects the
explanation that primary-selection count alone accounts for the effect. It does
not match which token positions survive the capacity filter, their raw gates,
or their gate-weighted gradient mass. In the principal top-1 setting, however,
preserving each expert's selected count also preserves its accepted count and
the aggregate dropped fraction because capacity is fixed. The intervention
therefore identifies an effect of the realized route beyond counts; it does not
isolate token identity from the accompanying gate and capacity-ranking changes.

\begin{figure*}[t]
\centering
\begin{subfigure}[t]{0.32\textwidth}
  \centering
  \resultpanel{figures/paper/fig3a_routing_mask_advantage.pdf}
  \caption{Magnitude-mask advantage by routing condition.}
  \label{fig:routing_pruning}
\end{subfigure}
\hfill
\begin{subfigure}[t]{0.32\textwidth}
  \centering
  \resultpanel{figures/paper/fig3b_route_mask_association.pdf}
  \caption{Final-route agreement versus 80\%-mask overlap.}
  \label{fig:route_mask}
\end{subfigure}
\hfill
\begin{subfigure}[t]{0.32\textwidth}
  \centering
  \resultpanel{figures/paper/fig3c_cross_init.pdf}
  \caption{Cross-initialization overlap for two target seeds.}
  \label{fig:cross_init}
\end{subfigure}
\caption{\textbf{Routing interventions and sparse-support formation on the
WikiText-103 subset.}
(a) Advantage is $100(L_{\rm random}-L_{\rm magnitude})/L_{\rm random}$;
shuffling primary IDs across token positions while preserving per-expert
selection counts nearly eliminates it. (b) All 15 unordered routing-condition
pairs are shown; colored labels mark comparisons with learned routing and the
exact learned-routing/replay point at $(1,1)$ is the positive control. This is
supporting association, not standalone causal evidence. (c) Tick marks are
the individual target seeds 17 and 29 and larger symbols are their means.
Cross-initialization replay reproduces the archived primary route exactly but
yields a different sparse outcome and a 6.54\% loss gap. Panels (a--b)
summarize three seeds with $\pm1$ sample standard deviation; no uncertainty
estimate is inferred from the two cross-initialization targets.}
\label{fig:routing_centerpiece}
\end{figure*}

Across routing-condition comparisons, agreement on the fixed final-checkpoint
validation routing set is associated with 80\%-sparse mask similarity
(\cref{fig:route_mask}).
The per-seed Pearson correlations between validation-route agreement and mask
Jaccard similarity are $r=0.7785\pm0.0275$ on WikiText-103 and
$r=0.7559\pm0.0122$ on TinyStories. This association supports the intervention
results but is not itself causal evidence, and coordinate overlap does not
imply functional equivalence.

Cross-initialization replay tests whether a recorded route is a portable sparse
blueprint. Imposing a source primary-ID trajectory on an independently
initialized target reproduces the archived IDs exactly, but increases
validation loss by
6.54\% relative to learned routing under the same target initialization and
produces a different sparse outcome (\cref{fig:cross_init}). Within
the evaluated setting, this is inconsistent with routing alone uniquely
determining the support. Together, the fixed-initialization interventions and
cross-initialization replay support the narrower conclusion that routed-token
history makes a causal contribution to support formation, while the resulting
support depends jointly on routing, initialization, and optimization.

\subsection{Robustness Across Architectures, Data, and Training Budget}
\label{sec:results_robustness}

Magnitude masks outperform matched random masks across models with 4, 8, and
16 experts; top-1 and top-2 routing; and 4 and 8 Transformer layers
(\cref{fig:robustness}). The same principal ordering is recorded for the
50k-story TinyStories subset, the WikiText-103 subset, a character-balanced
mixture, and a separate 1,024-token byte-ngram causal-control suite. Each
evaluated aggregate uses seeds 7, 17, and 29.

\begin{figure*}[t]
\centering
\begin{subfigure}[t]{0.32\textwidth}
  \centering
  \resultpanel{figures/paper/fig4a_architecture_robustness.pdf}
  \caption{80\%-mask advantage across 12 architectures.}
\end{subfigure}
\hfill
\begin{subfigure}[t]{0.32\textwidth}
  \centering
  \resultpanel{figures/paper/fig4b_dataset_robustness.pdf}
  \caption{Mask advantage across byte-tokenized datasets.}
\end{subfigure}
\hfill
\begin{subfigure}[t]{0.32\textwidth}
  \centering
  \resultpanel{figures/paper/fig4c_phase4_rewinds.pdf}
  \caption{Representative 80\%-sparse 10\% rewinds.}
\end{subfigure}
\caption{\textbf{Robustness summary.} In (a--b), advantage is the paired-seed
percentage reduction in pruning loss from replacing a random mask with the
magnitude mask; positive values favor magnitude. (a) Annotated means cover the
full $E\times k\times L$ grid, with the reference architecture outlined.
(b) The ordering holds at 50\% and 80\% sparsity on all three byte-tokenized
datasets. (c) Loss gaps relative to the paired dense model are shown for every
80\%-sparse control at the 10\% rewind point in three representative
2,500-step configurations; the dashed and dotted lines mark dense performance
and the 5\% tolerance. Panels (a--b) use three-seed paired means and panel (b)
shows sample standard deviations; panel (c) shows population standard
deviations of the paired seedwise percentage gaps, matching the follow-up's
population-SD convention.}
\label{fig:robustness}
\end{figure*}

A separate longer-budget suite trains the balanced-mixture 8-expert, top-1,
4-layer configuration for 10,000 steps, or approximately
$1.64\times10^8$ tokens per run. Its final checkpoints obscure the comparison:
normal routing reaches mean loss 1.5236 at step 10,000 but 1.3072 at the best
saved checkpoint (step 5,000 in every seed). At condition-specific best saved
checkpoints, normal routing again outperforms fixed-random,
random-every-step, and shuffled-usage routing. An 80\%-sparse normal-routing
mask extracted at the best saved checkpoint retrains from initialization to
$1.2984\pm0.0055$, versus $1.3072\pm0.0238$ for the best-saved dense model,
$1.3110\pm0.0070$ for a random mask, $1.3088\pm0.0074$ for random
reinitialization, and $1.5011\pm0.0129$ with randomized routing. The margins to
the mask controls are small, and this one-architecture result is a
checkpoint-selection case study rather than a factorial long-budget
replication.

\section{Discussion and Limitations}
\label{sec:discussion}

\paragraph{What the result means.}
Routing is not only an inference-time selector. During training it determines
the sequence of gradients received by each expert, and the experiments indicate
that this routed optimization leaves persistent sparse structure inside expert
parameters. The evidence goes beyond the expected observation that routing
changes weights: it identifies non-random expert-local supports, shows that a
subset remains trainable after rewinding, and demonstrates under controlled
interventions that realized routing history matters beyond logged
primary-selection counts. The shuffled control does not isolate token identity
from gate and capacity-ranking effects.

\paragraph{What the result does not mean.}
Sparse masks are a structural probe, not a complete characterization of expert
function or evidence of semantic specialization. The experiments do not show
that routing alone determines experts: cross-initialization replay gives the
opposite result. They also do not make route--mask correlation causal; the
causal claim comes from routing interventions at fixed training conditions.

\paragraph{Lottery-ticket interpretation.}
The reported 50\%-sparse supports that succeed after initialization rewind meet
the classical criterion in the original dataset suites. At 80\%, the original
and representative 2,500-step suites generally show early-rewind behavior; the
single 10,000-step best-checkpoint suite is an exception in which the learned
mask succeeds from initialization. Routing history conditions the process that
forms these supports; it need not be added to the formal definition of a
winning ticket. The supported picture is therefore support plus compatible
initialization, formed through routing-conditioned optimization.

\paragraph{Limitations.}
The experiments concern small decoder-only MoEs and corpus subsets rather than
production-scale systems. Principal masks use one-shot magnitude pruning; a
four-round representative IMP run improved on one-shot initialization rewind
at 80\% but was worse than the one-shot 10\%-rewind result; other pruning rules
may expose different structure. The transfer control replaces only expert 1's
mask with expert 0's mask, and learned-routing expert-local losses are evaluated
after rerouting rather than on frozen dense assignments. Route archives contain
pre-capacity primary expert IDs only, so the shuffled condition matches logged
selection and accepted counts under top-1, but not accepted-token identities,
raw gates, or gate-weighted gradient mass.
Principal results use three seeds and recorded standard deviations rather than
confidence intervals; follow-up aggregators historically mix sample and
population standard-deviation conventions. The long-budget experiment covers
one architecture, and the larger-validation follow-up is still modest. Finally,
coordinate-level mask overlap is sensitive to hidden-unit and expert-permutation
symmetries and should not be equated with functional equivalence.

\section{Conclusion}
\label{sec:conclusion}

Conventionally trained MoE experts contain non-random sparse supports, and the
reported 50\%-sparse supports retain optimization significance under
initialization rewinding. Routing interventions show that routed-token history
contributes to support formation beyond logged primary-selection counts, while
cross-initialization replay shows that routing does not uniquely determine the
outcome. Within the evaluated small-model regimes, the evidence therefore
supports expert-local sparse structure as a product of the interaction among
routing history, initialization, and optimization.

\section*{Reproducibility Statement}
The main paper defines the hypotheses, controls, normalization, and principal
training settings. Appendices~\ref{app:setup}--\ref{app:additional_controls}
record the implementation details, metrics, intervention construction,
statistical protocol, and completed extensions available in the repository.
The repository contains per-run resolved configurations, saved primary-route
archives, checkpoints, stored dataset text, and JSON result tables.
The publication panels are deterministically regenerated from those artifacts
by \path{scripts/generate_paper_figures.py};
\path{figures/paper/figure_manifest.json} records input and output hashes and
the uncertainty conventions. Dependencies
are not version pinned; upstream dataset revisions and exported-text checksums
are not recorded; and run summaries do not record a git revision, wall time, GPU
identity, or peak memory. These remain reproducibility gaps.

\section*{AI Use Statement}
A generative AI system assisted with this manuscript's restructuring, title and
prose drafting, LaTeX refactoring, and calibration of claims against existing
repository code and result artifacts. During this manuscript edit, it did not
generate new experimental data or launch new experiment runs. This statement
does not document possible AI use elsewhere in the research workflow, which the
authors must verify before submission. The authors retain responsibility for
all scientific claims, citations, numerical results, and final content.

\IfFileExists{references.bib}{%
  \IfFileExists{iclr2027_conference.bst}{%
    \bibliographystyle{iclr2027_conference}%
  }{%
    \bibliographystyle{plainnat}%
  }%
  \bibliography{references}%
}{%
  \typeout{references.bib is absent; bibliography omitted.}%
}

\appendix

\section{Full Experimental Setup}
\label{app:setup}

\subsection{Model and Routing}

The model is a pre-LayerNorm decoder with learned token and absolute-position
embeddings, causal multi-head self-attention, residual MoE blocks, a final
LayerNorm, and a bias-free language-model head tied to the token embedding.
Each expert is a two-layer feed-forward network
\begin{equation}
f_{\ell,e}(h;\theta_{\ell,e})
=
W_{\ell,e}^{(2)}
\operatorname{GELU}\!\left(W_{\ell,e}^{(1)}h+b_{\ell,e}^{(1)}\right)
+
b_{\ell,e}^{(2)},
\label{eq:app_expert}
\end{equation}
with dropout zero in every reported GPU run. For token representation
$h_{s,b,t}^{(\ell)}$, the bias-free learned router computes
\begin{equation}
p_{s,b,t}^{(\ell)}
=
\operatorname{softmax}\!\left(W_r^{(\ell)}h_{s,b,t}^{(\ell)}\right).
\label{eq:app_router}
\end{equation}
There is no routing noise. Let $S_{s,b,t}^{(\ell)}$ be the top-$k$ selected
expert set before capacity. For a batch of $B$ length-$T_{\rm seq}$ sequences,
the per-expert capacity is
\begin{equation}
C
=
\left\lceil
\alpha\frac{BT_{\rm seq}k}{E}
\right\rceil,
\qquad \alpha=1.25.
\label{eq:app_capacity}
\end{equation}
If an expert has more than $C$ selected assignments, the implementation accepts
the $C$ assignments with the largest current selected probabilities and drops
the rest. Writing $A_{s,b,t}^{(\ell)}\subseteq S_{s,b,t}^{(\ell)}$ for accepted
assignments, the MoE contribution is
\begin{equation}
y_{s,b,t}^{(\ell)}
=
\sum_{e\in A_{s,b,t}^{(\ell)}}
p_{s,b,t}^{(\ell)}(e)
f_{\ell,e}\!\left(h_{s,b,t}^{(\ell)};\theta_{\ell,e}\right).
\label{eq:app_moe_output}
\end{equation}
Selected probabilities are used raw: they are not renormalized over the top-$k$
set, so a top-1 gate is generally smaller than one. A dropped MoE contribution
is zero and the block residual remains.

The in-memory trace contains selected IDs, selected probabilities, pre-capacity
usage, and a dropped fraction. The compressed training and validation archives,
however, store only the first (primary) selected expert ID before capacity. They
do not store secondary top-$k$ IDs, accepted/drop status, or gate values. All
reported route-agreement and replay results therefore concern this archived
primary-ID trajectory.

Ignoring optimizer-state terms for notation, the expert update can be written
\begin{equation}
\theta_{\ell,e}^{s+1}
=
\theta_{\ell,e}^{s}
-
\eta_s
\sum_{x_i\in R_{\ell,e}(s)}
\nabla_{\theta_{\ell,e}}
\mathcal{L}(x_i;\Theta^s),
\label{eq:app_expert_update}
\end{equation}
where $R_{\ell,e}(s)$ is the accepted, gate-weighted routed set. Equal primary
selection counts need not give equal $R_{\ell,e}(s)$ or equal gradients: token
identity, gates, capacity, order, and the full model state all matter.

\subsection{Reference Configuration}

\begin{table}[t]
\centering
\scriptsize
\caption{Reference configuration for the principal pruning, routing, and
rewind experiments. AdamW $\beta$ and $\epsilon$ values are PyTorch defaults.}
\label{tab:app_reference_configuration}
\begin{tabular}{@{}ll@{}}
\toprule
Configuration & Value \\
\midrule
Architecture & 4 layers; 4 heads; $d_{\rm model}=256$ \\
Experts & 8 per layer; hidden width 1024; GELU \\
Router & Bias-free linear; top-1; no routing noise \\
Capacity / dropout & 1.25 / 0 \\
Vocabulary / positions & 256 UTF-8 byte IDs / 128 learned positions \\
Parameters & 17,981,952 \\
Sequence / batch & 128 tokens / 128 sequences \\
Updates / processed tokens & 2,500 / 40,960,000 \\
Optimizer & AdamW; constant learning rate $3\times10^{-4}$ \\
$\beta_1,\beta_2,\epsilon$ & $0.9,0.999,10^{-8}$ \\
Weight decay / gradient clip & 0.01 / global norm 1.0 \\
Balancing coefficient & 0.1 \\
Precision & FP16 autocast and gradient scaling; TF32 enabled \\
Evaluation cadence & Every 250 updates \\
Saved steps & 0, 25, 125, 250, 500, 1000, 2500 \\
Validation-batch cap & 12 (single-domain); 32 (mixture) \\
Random seeds & 7, 17, 29 \\
\bottomrule
\end{tabular}
\end{table}

Models are trained by next-token prediction with
\begin{equation}
\mathcal{L}
=
\mathcal{L}_{\mathrm{LM}}
+
\lambda_{\mathrm{bal}}\mathcal{L}_{\mathrm{bal}},
\qquad
\lambda_{\mathrm{bal}}=0.1.
\label{eq:app_training_objective}
\end{equation}
If $f_{\ell,e}$ is the fraction of pre-capacity selected slots assigned to
expert $e$ and $\bar p_{\ell,e}$ is its mean router probability, the
layer-averaged balancing term is
\begin{equation}
\mathcal{L}_{\mathrm{bal}}
=
\frac{1}{L}
\sum_{\ell=1}^{L}
E\sum_{e=1}^{E}
f_{\ell,e}\bar p_{\ell,e}.
\label{eq:app_balance}
\end{equation}
The selected fraction is detached in this auxiliary term.

\subsection{Data, Optimization, and Checkpoints}

The TinyStories subset contains the first 50,000 training stories and 2,000
validation stories exported from the local dataset. For WikiText-103-raw-v1,
the exporter selects the first 10,000 training and 1,000 validation source rows
from \texttt{Salesforce/wikitext}, then omits blank rows before joining the
retained text. The corresponding
stored UTF-8 files contain approximately 45.2 MB/1.62 MB (TinyStories train/val)
and 2.94 MB/0.247 MB (WikiText train/val).

The balanced corpus uses equal character budgets up to chunk rounding and
interleaves 16,384-character chunks from the two sources. Its recorded training
contributions are 2,916,352 TinyStories characters and 2,931,378 WikiText
characters; validation contributions are 245,760 and 246,050 characters. Thus
``balanced'' means approximately character-balanced and interleaved, not an
independent 50/50 draw at every batch.

Principal suites encode UTF-8 bytes directly as IDs 0--255. Text is divided
into contiguous, non-overlapping 128-token inputs with one-position-shifted
targets; the incomplete tail is discarded. Training blocks are shuffled with a
seeded generator and incomplete training batches are dropped; validation blocks
are ordered. Both dense and ticket trainers wrap the training loader in
\texttt{itertools.cycle}; after the first traversal, that iterator repeats the
same cached shuffled batch sequence rather than drawing a new epoch order. The
stored datasets yield 2,729 TinyStories, 179 WikiText, 357 byte-mixture, and 162
byte-ngram-mixture full training batches. Thus the 2,500-step runs traverse
approximately 0.92, 13.97, 7.00, and 15.43 such fixed orders respectively, and
the 10,000-step byte-mixture run traverses it about 28.01 times. The byte-ngram
robustness condition greedily longest-matches a
1,024-ID vocabulary learned from the first 750,000 training bytes, with maximum
ngram length 8, minimum frequency 3, and single-byte fallback. It produces
2,659,291 training and 223,471 validation tokens from 5,857,322 raw training
bytes.

Reference runs process $128\times128$ tokens per update for 2,500 updates,
approximately $4.10\times10^7$ tokens. Training uses AdamW at a constant
learning rate; there is no scheduler or warmup. CUDA runs use FP16 autocast,
gradient scaling, and TF32 matrix operations. Python, NumPy, CPU Torch, and CUDA
Torch RNGs are seeded, although deterministic-algorithm mode is not enabled.
Within a matched cell, architecture, initialization, deterministically generated
mini-batch order, optimizer, training budget, and evaluation set are fixed. The
routing regime is the manipulated variable except in the explicit
cross-initialization experiment.

Principal single-domain runs save checkpoints at the steps in
\cref{tab:app_reference_configuration}. A checkpoint contains model state,
configuration, step, validation loss, and optimizer state only when
\texttt{save\_optimizer} is enabled. Utilization JSONL logs, validation-route
NPZ files, and the optional training-route archive are separate artifacts.
Multi-domain and Phase-4 configs do not save optimizer state. Random pruning and
reinitialization controls use deterministic streams derived from the training
seed.

Run-specific suite configurations record the optimizer, learning rate, weight
decay, dropout, clipping, capacity factor, precision, checkpoint cadence,
routing mode, and preprocessing paths. The documented GPU environment is
PyTorch 2.5.1 with CUDA 12.4 on a 16 GB NVIDIA GeForce RTX 3080 Laptop GPU; a
reference benchmark records 1.1 GB peak allocation and approximately 25,100
tokens/s. Individual run summaries do not persist software versions, GPU
identity, wall time, peak memory, or a git revision, and project dependencies
are unpinned. The WikiText download does not pin a dataset revision, and the
TinyStories export does not pin an upstream revision or save an exported-text
checksum (although the local dataset state contains a Hugging Face
fingerprint). The
stored text files reproduce these runs locally, but regenerating them exactly
from upstream sources is not guaranteed. These are gaps for exact environment
reconstruction.

\section{Metrics and Statistical Protocol}
\label{app:metrics}

\subsection{Functional Metrics}

The primary metric is token-level cross entropy, summed over evaluated target
tokens and divided by their count, on a fixed held-out validation set. Reported
perplexity is $\exp(\min\{20,\mathcal L\})$. Because routing interventions can
change dense performance, each pruning
result is normalized to the dense model trained under the same condition. For
routing condition $c$, pruning method $a$, and sparsity $\rho$,
\begin{equation}
\Delta_{c,a}^{(\rho)}
=
\frac{
\mathcal{L}_{c,a}^{(\rho)}
-
\mathcal{L}_{c,\mathrm{dense}}
}{
\mathcal{L}_{c,\mathrm{dense}}
}.
\label{eq:app_relative_gap}
\end{equation}
The absolute learned-mask advantage over control $j$ is
\begin{equation}
G_{c,j}^{(\rho)}
=
\mathcal{L}_{c,j}^{(\rho)}
-
\mathcal{L}_{c,\mathrm{mag}}^{(\rho)}.
\label{eq:app_control_gap}
\end{equation}
The threshold $\varepsilon=0.05$ matches the generated Phase-4 full-loss flag
and is not a substitute for reporting continuous losses and control gaps.

For an evaluated model $q$, the logged expert-local diagnostic defines
\begin{equation}
\mathcal{V}_{q,\ell,e}
=
\{(x_i,y_i):e\text{ appears in model }q\text{'s selected top-}k\text{ IDs}\}.
\label{eq:app_local_set}
\end{equation}
Loss is restricted to the corresponding token subset. Under top-2 routing, a
token belongs to both selected experts' subsets. For learned-routing direct
pruning, each pruned model routes afresh; dense assignments are not frozen.
These expert-local values therefore combine functional change with possible
changes in the conditional token subset and are treated as secondary. Replay,
swap, and shuffle evaluation can use archived validation-ID overrides.

\subsection{Structural and Routing Metrics}

For binary masks $m_a$ and $m_b$, support similarity is Jaccard overlap:
\begin{equation}
J(m_a,m_b)
=
\frac{
|\operatorname{supp}(m_a)\cap\operatorname{supp}(m_b)|
}{
|\operatorname{supp}(m_a)\cup\operatorname{supp}(m_b)|
}.
\label{eq:app_jaccard}
\end{equation}
The standard aggregate helper sums intersections and unions over every common
expert weight tensor before dividing (a micro-aggregate), rather than averaging
layer--expert Jaccards. Functional-alignment analyses additionally report
Jaccard after expert and hidden-unit matching.

Route agreement is measured on the fixed final-checkpoint validation routing
set, not over the complete training trace. Let $r_i^{(\ell,c)}$ be the archived
primary expert ID.
\begin{equation}
R_{\mathrm{val}}(c,c')
=
\frac{1}{LN}
\sum_{\ell=1}^{L}
\sum_{i=1}^{N}
\mathbf{1}\!\left[r_i^{(\ell,c)}=r_i^{(\ell,c')}\right].
\label{eq:app_route_agreement}
\end{equation}
This is the fraction of validation token--layer positions with equal primary
IDs. Even for top-2 models it is not top-$k$ set overlap, because secondary IDs
are not archived. Its association with mask Jaccard is descriptive; causal
evidence comes from interventions on training-time routing.

For layer $\ell$, normalized expert-usage entropy is
\begin{equation}
\widetilde H_{\mathrm{use}}^{(\ell)}
=
-
\frac{1}{\log E}
\sum_{e=1}^{E}q_{\ell,e}\log q_{\ell,e},
\label{eq:app_usage_entropy}
\end{equation}
where $q_{\ell,e}$ is the pre-capacity selected-slot fraction assigned to expert
$e$. The analysis computes this entropy, coefficient of variation, max/min
ratio, and dead-expert count per logged step and layer. The run logs also record
router entropy, router margin, and dropped-assignment fraction; validation-route
archives support checkpoint-to-checkpoint primary-ID stability.

\subsection{Replication and Uncertainty}

The independent unit of replication is the training seed. Principal experiments
use seeds 7, 17, and 29 and report their mean and one recorded standard
deviation; error bars are not confidence intervals. Comparisons are paired by
seed. Experts, layers, tokens, sparsity levels, and checkpoints are not treated
as independent replicates. The original multi-seed, load-balance, architecture,
multi-domain, and swap aggregators use sample standard deviation. Several
follow-up rewind, IMP, and long-budget aggregators use population standard
deviation. Tables below preserve the recorded values and identify two-target
cross-initialization summaries separately; no inferential tests are claimed.

For validation-route-agreement/mask-similarity analyses, Pearson correlation is
computed separately for each seed over all 15 unordered pairs among the six
routing conditions using final-checkpoint validation routes and 80\% final
magnitude masks, then summarized across seeds. This
construction includes
the exact normal--replay positive-control point $(R,J)=(1,1)$, which can have
substantial leverage, and the 15 within-seed pairs are not independent. The
reported values are
$r=0.7785\pm0.0275$ on the WikiText-103 subset and
$r=0.7559\pm0.0122$ on the TinyStories subset. These correlations do not replace
intervention-based evidence and do not imply that mask overlap is functional
equivalence.

\section{Full Direct-Pruning Protocol and Results}
\label{app:direct_pruning}

\subsection{Mask Extraction}

Only expert weight matrices are pruned. Router, attention, embedding,
normalization, output, and expert-bias parameters remain dense. Let
\begin{equation}
\vartheta_{\ell,e,T}
=
\operatorname{vec}
\left(
W_{\ell,e,T}^{(1)},W_{\ell,e,T}^{(2)}
\right)
\label{eq:app_vectorized_expert}
\end{equation}
be the concatenated prunable weights of expert $(\ell,e)$. Expert-local
magnitude pruning retains an exact-count largest-$1-\rho$ subset by absolute
value. Let $n_{\ell,e}=\dim(\vartheta_{\ell,e,T})$. With
\begin{equation}
K_{\ell,e}^{(\rho)}
=
\max\!\left\{1,\operatorname{round}\!\left((1-\rho)n_{\ell,e}\right)\right\},
\end{equation}
the implemented mask is
\begin{equation}
m_{\ell,e}^{(\rho)}(j)
=
\mathbf{1}\!\left[
j\in\operatorname{TopK}_{K_{\ell,e}^{(\rho)}}
\!\left(|\vartheta_{\ell,e,T}|\right)
\right].
\label{eq:app_mask_threshold}
\end{equation}
Selection is joint across both weight matrices within an expert. The retained
count is therefore exact, and \texttt{torch.topk(sorted=False)} selects those
coordinates. There is no
explicit deterministic tie-break rule beyond the backend top-$k$
implementation. The joint expert-local rank gives each expert the requested
overall sparsity even when tensor or expert scales differ.

The original single-domain direct-pruning grid is
\begin{equation}
\rho\in\{0.50,0.70,0.80,0.90,0.95\},
\label{eq:app_sparsity_grid}
\end{equation}
while the architecture, multi-domain, and rewind extensions use 50\% and 80\%.
Principal claims emphasize those two sparsities.

\paragraph{Controls.}
A random mask samples surviving coordinates uniformly within each prunable
tensor while exactly matching that tensor's learned-mask count. The recorded
``other expert'' control is limited: in every layer it copies expert 0's two
masks onto expert 1 and leaves every other expert's learned mask unchanged. It
is not an all-expert cyclic transfer. Random reinitialization calls the PyTorch
default reset routine on every expert Linear module, thereby resampling expert
weights and biases, and then reapplies the learned weight mask. In direct
pruning, this separates learned final values from sparse topology; it is not an
initialization-rewind test. All direct-pruning masks are evaluated without
additional optimization, and learned-routing models may reroute after pruning.

\subsection{Extended Outcome Summary}

At the principal 50\% and 80\% sparsities, learned masks outperform random,
limited-transfer, and random-reinitialization controls on both original data
subsets. At 50\%, the learned-mask increase is 3.1\% on WikiText versus 140.6\%
for a random mask and 545.0\% for random reinitialization; the TinyStories
learned-mask increase is 4.6\%. At 80\%, learned masks retain an advantage but
direct pruning is far from dense. At 90--95\%, all methods approach destructive
saturation and the ordering is not uniform: for example, WikiText 95\%
magnitude loss is slightly above both random and limited-transfer loss.
\Cref{tab:app_direct_pruning_full} gives every recorded three-seed mean and
sample standard deviation in the original direct-pruning grid.

\begin{table*}[t]
\centering
\scriptsize
\caption{Full normal-routing direct-pruning grid. Entries are validation loss,
mean $\pm$ sample SD over seeds 7, 17, and 29. ``0$\to$1 transfer'' changes only
expert 1's mask in each layer.}
\label{tab:app_direct_pruning_full}
\resizebox{\textwidth}{!}{%
\begin{tabular}{@{}llcccc@{}}
\toprule
Dataset & Sparsity & Magnitude & Random mask & 0$\to$1 transfer & Random reinit \\
\midrule
TinyStories & 50\% & $1.0676\pm0.0168$ & $3.2700\pm0.0873$ & $1.3590\pm0.0282$ & $5.5725\pm0.1793$ \\
 & 70\% & $1.7285\pm0.0222$ & $4.7171\pm0.1398$ & $2.3025\pm0.0692$ & $5.5820\pm0.2175$ \\
 & 80\% & $3.5043\pm0.3249$ & $5.1919\pm0.1427$ & $3.9594\pm0.4261$ & $5.5691\pm0.2084$ \\
 & 90\% & $5.2208\pm0.3144$ & $5.4524\pm0.1795$ & $5.3101\pm0.3236$ & $5.5843\pm0.2109$ \\
 & 95\% & $5.5092\pm0.2397$ & $5.4989\pm0.1897$ & $5.5150\pm0.2226$ & $5.5780\pm0.1950$ \\
\midrule
WikiText-103 & 50\% & $1.7388\pm0.0205$ & $4.0564\pm0.4725$ & $1.9418\pm0.1086$ & $10.8751\pm4.1653$ \\
 & 70\% & $2.4129\pm0.1115$ & $7.4999\pm2.0359$ & $2.7653\pm0.1519$ & $10.9116\pm4.1331$ \\
 & 80\% & $4.5586\pm0.5000$ & $9.3660\pm3.2425$ & $4.9161\pm0.5673$ & $10.9424\pm4.1958$ \\
 & 90\% & $9.4493\pm1.9800$ & $10.4422\pm3.9551$ & $9.3411\pm2.1794$ & $10.9860\pm4.2031$ \\
 & 95\% & $10.8825\pm3.4328$ & $10.6401\pm4.0744$ & $10.7710\pm3.5516$ & $10.9704\pm4.1709$ \\
\bottomrule
\end{tabular}}
\end{table*}

These results establish non-random final-weight structure at the principal
sparsities. The stronger optimization-significance claim depends on
\cref{app:rewind}.

\section{Full Rewind-and-Retrain Suite}
\label{app:rewind}

Let $M^{(\rho)}$ be one on all unpruned parameters and equal to the expert masks
on prunable weights. At rewind point $t_0$, the complete model is restored and
masked:
\begin{equation}
\widetilde\Theta_{t_0}^{(\rho)}
=
M^{(\rho)}\odot\Theta_{t_0}.
\label{eq:app_rewind_state}
\end{equation}
Embeddings, attention, normalization, routers, expert biases, and surviving
expert weights are therefore rewound to the same checkpoint. During retraining,
pruned expert coordinates remain zero:
\begin{equation}
\nabla_\Theta\mathcal{L}
\leftarrow
M^{(\rho)}\odot\nabla_\Theta\mathcal{L}.
\label{eq:app_gradient_mask}
\end{equation}
Gradient hooks apply this operation, and the weight mask is reapplied after
every optimizer step so masked coordinates cannot regrow through optimizer
updates or weight decay.
A fresh AdamW optimizer uses the original constant learning rate and other
hyperparameters. Every condition rebuilds the seeded loader at the beginning of
its deterministic order and retrains for the full original $T$ updates,
regardless of $t_0$; it does not consume only the suffix or resume optimizer
state. Surviving parameters rewound to $t_0>0$ therefore receive $T$ additional
updates, so these early-rewind comparisons are full-budget restart experiments,
not compute-matched $T-t_0$ continuations.

Principal rewind fractions are
\begin{equation}
t_0/T\in\{0,0.01,0.05,0.10\},
\label{eq:app_rewind_points}
\end{equation}
with 0.25 additionally evaluated in the reduced long-budget suite. The principal
fractions map to checkpoints 0, 25, 125, and 250. Experiments evaluate
50\% and 80\% sparsity and compare (i) the learned mask, (ii) a matched random
mask, (iii) the learned mask with randomly reinitialized expert parameters,
(iv) the learned mask retrained under random-every-step routing, and (v) the
corresponding dense reference where a lottery-ticket classification is made.

The discovery trace $\mathcal{A}^{\mathrm{disc}}$, which produces the mask, is
distinguished from the retraining trace $\mathcal{A}^{\mathrm{ret}}$. The
default retraining condition matches the discovery routing regime; the
randomized-routing control changes the sparse retraining trajectory. Routing
conditions ticket formation and retraining, but it is not added to the formal
definition of a winning ticket.

\Cref{tab:app_rewind_original} enumerates the original two-dataset suite. At
50\% sparsity, learned supports retrain from initialization on both subsets and
outperform the sparse controls. At 80\%, initialization rewind is worse than
dense, while a 10\% checkpoint plus a full retraining budget reaches
$1.0413\pm0.0138$ on TinyStories and $1.6882\pm0.0297$ on WikiText, within
2.03\% and 0.13\% of their respective dense means. Cells without a matched
dense reference are described by continuous loss and control gaps rather than
assigned a binary label.

\begin{table*}[t]
\centering
\scriptsize
\caption{Complete original rewind aggregates. Entries are validation loss,
mean $\pm$ sample SD over seeds 7, 17, and 29. Dense reference means are
$1.0206\pm0.0096$ (TinyStories) and $1.6859\pm0.0160$ (WikiText-103). Every
row uses a fresh optimizer and a full 2,500-update retraining budget.}
\label{tab:app_rewind_original}
\resizebox{\textwidth}{!}{%
\begin{tabular}{@{}llrcccc@{}}
\toprule
Dataset & Sparsity & Rewind & Learned mask & Random mask & Random reinit & Randomized routing \\
\midrule
TinyStories & 50\% & 0\% & $1.0365\pm0.0068$ & $1.0915\pm0.0109$ & $1.0785\pm0.0055$ & $1.5506\pm0.0122$ \\
 & & 1\% & $1.0212\pm0.0076$ & $1.0955\pm0.0096$ & $1.1126\pm0.0230$ & $1.5482\pm0.0381$ \\
 & & 5\% & $0.9734\pm0.0237$ & $1.0474\pm0.0142$ & $1.0460\pm0.0092$ & $1.4542\pm0.0149$ \\
 & & 10\% & $0.9544\pm0.0070$ & $1.0353\pm0.0149$ & $1.0356\pm0.0115$ & $1.4321\pm0.0324$ \\
 & 80\% & 0\% & $1.1492\pm0.0034$ & $1.2488\pm0.0096$ & $1.2369\pm0.0256$ & $1.5998\pm0.0168$ \\
 & & 1\% & $1.1629\pm0.0101$ & $1.2510\pm0.0150$ & $1.2233\pm0.0141$ & $1.5821\pm0.0218$ \\
 & & 5\% & $1.0780\pm0.0240$ & $1.1748\pm0.0013$ & $1.1656\pm0.0030$ & $1.4899\pm0.0319$ \\
 & & 10\% & $1.0413\pm0.0138$ & $1.1606\pm0.0035$ & $1.1534\pm0.0115$ & $1.4579\pm0.0209$ \\
\midrule
WikiText-103 & 50\% & 0\% & $1.6593\pm0.0103$ & $1.7330\pm0.0334$ & $1.7558\pm0.0190$ & $2.1470\pm0.0109$ \\
 & & 1\% & $1.6441\pm0.0235$ & $1.7348\pm0.0442$ & $1.7397\pm0.0449$ & $2.1388\pm0.0188$ \\
 & & 5\% & $1.6062\pm0.0103$ & $1.6891\pm0.0530$ & $1.6966\pm0.0653$ & $2.0953\pm0.0144$ \\
 & & 10\% & $1.5999\pm0.0103$ & $1.6590\pm0.0263$ & $1.6817\pm0.0615$ & $2.0693\pm0.0043$ \\
 & 80\% & 0\% & $1.8399\pm0.0254$ & $1.9203\pm0.0193$ & $1.9168\pm0.0149$ & $2.1743\pm0.0157$ \\
 & & 1\% & $1.8294\pm0.0335$ & $1.9234\pm0.0372$ & $1.8814\pm0.0113$ & $2.1592\pm0.0080$ \\
 & & 5\% & $1.7162\pm0.0356$ & $1.8490\pm0.0486$ & $1.8353\pm0.0442$ & $2.1154\pm0.0054$ \\
 & & 10\% & $1.6882\pm0.0297$ & $1.8301\pm0.0540$ & $1.8245\pm0.0515$ & $2.0905\pm0.0143$ \\
\bottomrule
\end{tabular}}
\end{table*}

\section{Routing Interventions}
\label{app:interventions}

\subsection{Trace Construction}

For every layer, the saved assignment trace is the primary pre-capacity ID
$r_{s,b,t}^{(\ell)}$. Equivalently,
\begin{equation}
\mathcal{A}^{(\ell)}
=
\left(a_{s,b,t,e}^{(\ell)}\right)_{s,b,t,e},
\qquad
a_{s,b,t,e}^{(\ell)}
=
\mathbf{1}\!\left[e=r_{s,b,t}^{(\ell)}\right].
\label{eq:app_assignment_trace}
\end{equation}
The archive stores this ID tensor for each training step and layer. It stores no
gate values, secondary top-$k$ IDs, or capacity acceptance. Replay is indexed by
training step, batch index, and sequence position, so source and target require
the same tensor shapes and ordered token stream. Archived primary utilization is
\begin{equation}
U_{\ell,e}
=
\sum_{s,b,t}a_{s,b,t,e}^{(\ell)}.
\label{eq:app_utilization}
\end{equation}

\subsection{Intervention Definitions}

\paragraph{Learned routing.}
The router and experts are jointly optimized under the language-model and
load-balancing objectives.

\paragraph{Same-initialization replay.}
Primary IDs from a source run are imposed on a target with the source
initialization and ordered data. The target router recomputes its current
softmax probability for the imposed ID. This positive control tests ID-replay
fidelity. Empirically it reproduces archived IDs and extracted masks with
agreement 1.0/1.0 and exactly matches the normal dense loss in all three
original-suite seeds.

\paragraph{Random-every-step routing.}
In the principal top-1 experiments, a vector formed by
$0,1,\ldots,E-1$ modulo $E$ is deterministically permuted across the $BT$
positions using seed $s_0+1009s+9176\ell$. Primary selection counts are
therefore balanced to within one assignment per expert in each layer and step.
Capacity and current gate weights are then applied normally. In top-2 code paths
this override fixes only the primary ID; the remaining selected ID is recomputed
from current router probabilities.

\paragraph{Frozen random router.}
The router is randomly initialized and frozen before training. Its parameter
partition is fixed, but assignments may change as hidden representations
evolve. This is not a fixed assignment table; it tests stable but untrained
routing geometry.

\paragraph{Shuffled-usage routing.}
For each layer and training step, the source primary-ID tensor is deterministically
permuted across token positions using the same seed formula. Consequently,
\begin{equation}
U_{\ell,e,s}^{\mathrm{shuffled}}
=
U_{\ell,e,s}^{\mathrm{source}}.
\label{eq:app_shuffled_usage}
\end{equation}
This exactly matches archived primary-selection counts at the layer-and-step
level, not merely in aggregate, while generally changing the ordered token
history $\mathcal H_{\ell,e}$. In the principal top-1 setting, fixed capacity
therefore also preserves each expert's accepted count and the total dropped
fraction. The current gates and capacity ranking are nevertheless recomputed,
so the accepted-token identities, acceptance mask over positions, and
gate-weighted gradients need not match. Raw top-1 gates are not unit-valued.

\paragraph{Swapped and cyclic routing.}
A fixed expert permutation $\pi$ is applied to a recorded trace:
\begin{equation}
a_{s,b,t,\pi(e)}^{(\ell),\mathrm{swap}}
=
a_{s,b,t,e}^{(\ell),\mathrm{source}}.
\label{eq:app_swap}
\end{equation}
Only archived IDs are permuted; the current router probability of the substituted
ID supplies the gate. Pairwise swaps are used in the principal suite, while the
multi-domain extension includes several global, layer-specific, and cyclic ID
permutations.

Except for cross-initialization replay, interventions use the same model
initialization and ordered data. Pruning results are normalized to the dense
model from the same routing condition. Only the frozen-random condition freezes
router parameters. In random-every-step, replay, shuffled, swapped, and
cross-initialization replay, the router projection remains trainable; an
imposed primary ID gathers that expert's current softmax probability, and
accepted assignments receive language-model gradients through that gate. The
router also remains coupled to the auxiliary balancing objective.

\subsection{Fidelity and Secondary Results}

\begin{figure*}[t]
\centering
\begin{subfigure}[t]{0.48\textwidth}
  \centering
  \resultpanel{figures/paper/figA1_dense_routing.pdf}
  \caption{Dense performance under routing interventions.}
\end{subfigure}
\hfill
\begin{subfigure}[t]{0.48\textwidth}
  \centering
  \resultpanel{figures/paper/figA2_routing_pruning_curves.pdf}
  \caption{Dense-normalized magnitude-pruning curves.}
\end{subfigure}
\caption{Secondary routing controls. Replay reproduces the learned-routing
baseline within the reported variability. Random-every-step and shuffled-usage
routing degrade dense performance and reduce the learned-mask advantage;
condition-specific dense normalization separates these effects. Points are
three-seed means, error bars are sample standard deviations, and the pruning
axis in (b) is base-2 logarithmic.}
\label{fig:app_routing_controls}
\end{figure*}

\Cref{fig:app_routing_controls,tab:app_dense_causal} summarize the recorded
dense intervention aggregates and their pruning consequences.
The shuffled-usage result is the critical selected-count control: primary-ID
counts are preserved while dense performance and the learned-mask advantage are
substantially weakened. Because gate and acceptance information is absent, it
is not an exact realized-update control: accepted counts are matched under
top-1, but accepted-token identities and gate-weighted gradients are not. The
validation-route/mask association
in \cref{sec:results_routing} is supporting evidence only.

\begin{table*}[t]
\centering
\scriptsize
\caption{Dense routing-condition validation loss, mean $\pm$ sample SD. Dashes
denote conditions not included in that aggregate. All rows use seeds 7, 17, and
29. Losses are comparable within, not across, tokenizer rows.}
\label{tab:app_dense_causal}
\resizebox{\textwidth}{!}{%
\begin{tabular}{@{}lcccccc@{}}
\toprule
Dataset / tokenizer & Learned & Fixed random & Random/step & Replay & Swap 0$\leftrightarrow$1 & Shuffled usage \\
\midrule
TinyStories / byte & $1.0206\pm0.0096$ & $1.1105\pm0.0184$ & $1.4586\pm0.0153$ & $1.0206\pm0.0096$ & $1.0491\pm0.0111$ & $1.5316\pm0.0151$ \\
WikiText-103 / byte & $1.6859\pm0.0160$ & $1.7386\pm0.0252$ & $2.1121\pm0.0080$ & $1.6859\pm0.0160$ & $1.7081\pm0.0151$ & $2.1337\pm0.0191$ \\
Balanced mixture / byte & $1.4801\pm0.0425$ & $1.5810\pm0.0207$ & $1.9564\pm0.0120$ & -- & -- & $1.9791\pm0.0189$ \\
Balanced mixture / byte-ngram & $3.2753\pm0.0307$ & $3.5037\pm0.0876$ & $4.3504\pm0.0323$ & -- & -- & $4.3574\pm0.0156$ \\
\bottomrule
\end{tabular}}
\end{table*}

The corresponding single-domain pruning results are reported numerically in
\cref{tab:app_single_domain_causal_pruning}. Under learned and frozen-random
routing, mask identity has a large effect. Under shuffled routing, magnitude,
random, limited-transfer, and reinitialized masks cluster much more closely.

\begin{table*}[t]
\centering
\scriptsize
\caption{Single-domain direct pruning by routing condition, mean loss over
seeds 7, 17, and 29. Sample SD is omitted for table width; transfer is the
limited expert-0-to-expert-1 control.}
\label{tab:app_single_domain_causal_pruning}
\resizebox{\textwidth}{!}{%
\begin{tabular}{@{}llrrrrrrrrr@{}}
\toprule
Dataset & Routing & Dense & 50\% mag. & 50\% random & 50\% transfer & 50\% reinit & 80\% mag. & 80\% random & 80\% transfer & 80\% reinit \\
\midrule
TinyStories & Learned & 1.0206 & 1.0676 & 3.2700 & 1.3590 & 5.5725 & 3.5043 & 5.1919 & 3.9594 & 5.5691 \\
 & Fixed random & 1.1105 & 1.1551 & 2.6055 & 1.3393 & 4.8932 & 4.1013 & 4.4096 & 3.3602 & 4.8813 \\
 & Random/step & 1.4586 & 1.4844 & 2.1112 & 1.5522 & 2.8181 & 2.0078 & 2.7004 & 2.0925 & 2.8181 \\
 & Shuffled usage & 1.5316 & 1.5349 & 1.5969 & 1.5423 & 1.6377 & 1.5743 & 1.6318 & 1.5825 & 1.6376 \\
\midrule
WikiText-103 & Learned & 1.6859 & 1.7388 & 4.0564 & 1.9418 & 10.8751 & 4.5586 & 9.3660 & 4.9161 & 10.9424 \\
 & Fixed random & 1.7386 & 1.7803 & 3.1131 & 1.9566 & 7.2169 & 4.0933 & 6.1239 & 4.3324 & 7.2072 \\
 & Random/step & 2.1121 & 2.1257 & 2.4422 & 2.1562 & 2.9099 & 2.3651 & 2.8214 & 2.4090 & 2.9096 \\
 & Shuffled usage & 2.1337 & 2.1352 & 2.1645 & 2.1368 & 2.1863 & 2.1529 & 2.1835 & 2.1547 & 2.1864 \\
\bottomrule
\end{tabular}}
\end{table*}

In the byte mixture, shuffled usage is 33.72\% worse than learned routing; in
the byte-ngram suite it is 33.04\% worse. The stored count check compares 3,232
layer/step/expert records per seed and reports zero mismatches for all three
seeds in both suites. The corresponding route-agreement/mask-Jaccard
correlations are much smaller than in the original single-domain suites:
$0.1804\pm0.0424$ for the byte mixture and $0.1237\pm0.1610$ for byte-ngram,
showing that the simple pairwise association is not universal.

\begin{table}[t]
\centering
\scriptsize
\caption{Balanced-mixture ID-permutation interventions. Loss, route agreement
with normal, and 80\% mask Jaccard are mean $\pm$ sample SD over three seeds.}
\label{tab:app_swap_controls}
\resizebox{\linewidth}{!}{%
\begin{tabular}{@{}lcccc@{}}
\toprule
Condition & Loss & $\Delta$ vs normal & Primary-route agreement & Mask Jaccard \\
\midrule
Global 0$\leftrightarrow$1 & $1.5111\pm0.0248$ & $+2.10\%$ & $0.7461\pm0.0084$ & $0.5679\pm0.0561$ \\
Global 0$\leftrightarrow$4 & $1.5068\pm0.0115$ & $+1.80\%$ & $0.7453\pm0.0115$ & $0.5405\pm0.0375$ \\
Global 2$\leftrightarrow$6 & $1.5027\pm0.0149$ & $+1.53\%$ & $0.7509\pm0.0224$ & $0.5616\pm0.0278$ \\
Layer 0, 0$\leftrightarrow$1 & $1.5026\pm0.0156$ & $+1.52\%$ & $0.9317\pm0.0080$ & $0.5869\pm0.0462$ \\
Layer 3, 0$\leftrightarrow$1 & $1.4897\pm0.0239$ & $+0.65\%$ & $0.9378\pm0.0081$ & $0.6822\pm0.0484$ \\
Global cyclic $+1$ & $1.5358\pm0.0099$ & $+3.77\%$ & $0.0000\pm0.0000$ & $0.4267\pm0.0160$ \\
\bottomrule
\end{tabular}}
\end{table}

Swap comparisons must be interpreted with the applied permutation in mind:
unaffected IDs mechanically inflate aggregate route agreement in the pairwise
and layer-specific cases.

\section{Cross-Initialization Replay}
\label{app:cross_initialization}

Same-initialization replay validates the ID-override mechanism but does not
separate routing from initialization. The completed WikiText experiment uses
source seed 7 and independently initialized target seeds 17 and 29. Every target
uses the source seed's data order. Learned target routing and imposed source
routing therefore share target initialization, architecture, data order,
optimizer settings, and budget; they differ in the archived primary IDs imposed
during training.

\begin{table}[t]
\centering
\scriptsize
\caption{Cross-initialization replay performance and primary-route fidelity.
The source seed is 7; the two target rows are not a three-seed aggregate.}
\label{tab:app_cross_init_performance}
\resizebox{\linewidth}{!}{%
\begin{tabular}{@{}rccccc@{}}
\toprule
Target & Source loss & Target learned & Cross-init replay & Source--learned route & Source--replay route \\
\midrule
17 & 1.6817 & 1.7130 & 1.8093 & 0.1083 & 1.0000 \\
29 & 1.6817 & 1.6418 & 1.7648 & 0.1137 & 1.0000 \\
\bottomrule
\end{tabular}}
\end{table}

The imposed archive reproduces the source primary IDs exactly, while replay
loss is 6.54\% above matched-data learned routing when averaged over the two
targets. The within-target learned and replay masks also differ substantially.

\begin{table}[t]
\centering
\scriptsize
\caption{Cross-initialization mask Jaccard. ``Same init'' compares learned and
replay masks inside the target coordinate system.}
\label{tab:app_cross_init_masks}
\begin{tabular}{@{}rcrrr@{}}
\toprule
Target & Sparsity & Source--learned & Source--replay & Same-init learned--replay \\
\midrule
17 & 50\% & 0.3752 & 0.3732 & 0.5897 \\
17 & 80\% & 0.1945 & 0.1872 & 0.4741 \\
29 & 50\% & 0.3743 & 0.3734 & 0.5833 \\
29 & 80\% & 0.1916 & 0.1876 & 0.4627 \\
\midrule
Mean & 50\% & 0.3747 & 0.3733 & 0.5865 \\
Mean & 80\% & 0.1931 & 0.1874 & 0.4684 \\
\bottomrule
\end{tabular}
\end{table}

To test coordinate symmetries directly, the completed functional-alignment
analysis evaluates both models on the same 512 validation-token positions,
matches experts within each layer by output linear CKA, and then matches hidden
units by absolute activation correlation. As
\cref{tab:app_functional_alignment} shows,
alignment changes Jaccard by only 0.0001--0.0007 in the aggregate. This reduces,
but does not eliminate, the concern that raw overlap is dominated by a trivial
expert or neuron permutation.

\begin{table}[t]
\centering
\scriptsize
\caption{Mean cross-initialization functional alignment over targets 17 and
29.}
\label{tab:app_functional_alignment}
\begin{tabular}{@{}lcrrrr@{}}
\toprule
Condition & Sparsity & Raw $J$ & Aligned $J$ & Expert CKA & Abs. neuron corr. \\
\midrule
Matched-data learned & 50\% & 0.3747 & 0.3750 & 0.5125 & 0.4041 \\
Matched-data learned & 80\% & 0.1931 & 0.1938 & 0.5125 & 0.4041 \\
Cross-init replay & 50\% & 0.3733 & 0.3735 & 0.4146 & 0.3989 \\
Cross-init replay & 80\% & 0.1874 & 0.1875 & 0.4146 & 0.3989 \\
\bottomrule
\end{tabular}
\end{table}

The same two-target aggregate also contains the direct-pruning comparison in
\cref{tab:app_cross_init_pruning}. Cross-init replay has worse dense and 50\%
magnitude loss than matched-data routing, but its 80\% mask is less fragile in
absolute terms. Both conditions retain a learned-mask advantage over random
and limited-transfer masks.

\begin{table}[t]
\centering
\scriptsize
\caption{Cross-initialization direct pruning, mean over target seeds 17 and
29. Transfer is the limited expert-0-to-expert-1 control.}
\label{tab:app_cross_init_pruning}
\begin{tabular}{@{}lrrrr@{}}
\toprule
Condition & Sparsity & Magnitude & Random mask & Transfer \\
\midrule
Matched-data learned & 50\% & 1.7256 & 3.7363 & 1.9281 \\
Matched-data learned & 80\% & 4.1736 & 7.5181 & 4.4492 \\
Cross-init replay & 50\% & 1.8151 & 2.6971 & 1.9772 \\
Cross-init replay & 80\% & 2.5127 & 3.9254 & 2.6936 \\
\bottomrule
\end{tabular}
\end{table}

The follow-up rewind suite uses the same two target seeds. Every learned-mask
curve is best at the tested 10\% rewind; \cref{tab:app_cross_init_rewind}
includes all matched controls at that point. Relative to each condition's own
dense mean, the cross-init gaps are $-4.78\%$ and $+0.59\%$ at 50\% and 80\%,
while the matched-data gaps are $-4.95\%$ and $+0.42\%$. The complete artifact
contains all 64 aggregate cells across four rewind points.

\begin{table*}[t]
\centering
\scriptsize
\caption{Cross-initialization full-budget retraining at the 10\% rewind. Values
are two-target means $\pm$ population SD.}
\label{tab:app_cross_init_rewind}
\resizebox{\textwidth}{!}{%
\begin{tabular}{@{}llccccc@{}}
\toprule
Condition & Sparsity & Dense & Learned mask & Random mask & Random reinit & Randomized routing \\
\midrule
Matched-data learned & 50\% & $1.6774\pm0.0356$ & $1.5944\pm0.0235$ & $1.6597\pm0.0131$ & $1.6736\pm0.0613$ & $2.0703\pm0.0092$ \\
Matched-data learned & 80\% & $1.6774\pm0.0356$ & $1.6844\pm0.0393$ & $1.8304\pm0.0579$ & $1.8107\pm0.0660$ & $2.0918\pm0.0063$ \\
Cross-init replay & 50\% & $1.7871\pm0.0223$ & $1.7018\pm0.0322$ & $1.8282\pm0.0340$ & $1.8361\pm0.0435$ & $2.0780\pm0.0043$ \\
Cross-init replay & 80\% & $1.7871\pm0.0223$ & $1.7977\pm0.0370$ & $1.9651\pm0.0082$ & $1.9825\pm0.0060$ & $2.1006\pm0.0040$ \\
\bottomrule
\end{tabular}}
\end{table*}

Thus a foreign ID history can form a rewindable mask, but it does not match the
route that co-adapts with the target initialization. The bounded conclusion is
that routing changes support within an initialization, while the recorded route
is not a portable coordinate-level blueprint.

\section{Architecture and Dataset Robustness}
\label{app:robustness}

Evaluated architecture cells draw from
\begin{equation}
E\in\{4,8,16\},
\qquad
k\in\{1,2\},
\qquad
L\in\{4,8\},
\label{eq:app_architecture_grid}
\end{equation}
using seeds 7, 17, and 29. The full factorial contains 36 WikiText training
runs. \Cref{tab:app_architecture_grid} records every aggregate architecture
cell. Magnitude masks beat random masks at 50\% and 80\% in all 12 cells, but
one-shot 80\% pruning is often catastrophic. No aggregate cell has a dead
expert; mean selected-route entropy spans 0.9846--0.9952, and mean dropped
fraction spans 0--0.0156.

\begin{table*}[t]
\centering
\scriptsize
\caption{Full WikiText architecture grid, mean over seeds 7, 17, and 29.
Dense uncertainty is sample SD. $H$ and dropped-assignment fraction $D$ are
final logged-step values averaged over layers within each run and then seeds.
M/R denote magnitude/random-mask loss.}
\label{tab:app_architecture_grid}
\resizebox{\textwidth}{!}{%
\begin{tabular}{@{}rrrrrcccccc@{}}
\toprule
$E$ & $k$ & $L$ & Params (M) & Dense loss & $H$ & $D$ & 50\% M & 50\% R & 80\% M & 80\% R \\
\midrule
4 & 1 & 4 & 9.57 & $1.7127\pm0.0304$ & 0.9882 & 0.0000 & 1.7954 & 5.0349 & 8.8331 & 9.9675 \\
4 & 1 & 8 & 19.04 & $1.5808\pm0.0110$ & 0.9846 & 0.0065 & 1.6445 & 5.2285 & 6.2745 & 8.7848 \\
4 & 2 & 4 & 9.57 & $1.6716\pm0.0183$ & 0.9913 & 0.0074 & 1.8004 & 12.3798 & 17.8248 & 28.3398 \\
4 & 2 & 8 & 19.04 & $1.5606\pm0.0102$ & 0.9952 & 0.0013 & 1.6357 & 7.8873 & 9.3844 & 14.5287 \\
8 & 1 & 4 & 17.98 & $1.6859\pm0.0160$ & 0.9932 & 0.0057 & 1.7388 & 4.0564 & 4.5586 & 9.3660 \\
8 & 1 & 8 & 35.87 & $1.5721\pm0.0070$ & 0.9913 & 0.0091 & 1.6115 & 4.3213 & 4.6839 & 7.5686 \\
8 & 2 & 4 & 17.98 & $1.6499\pm0.0154$ & 0.9913 & 0.0110 & 1.7407 & 8.0455 & 12.9197 & 16.5732 \\
8 & 2 & 8 & 35.87 & $1.5608\pm0.0061$ & 0.9927 & 0.0089 & 1.6042 & 6.1661 & 6.7442 & 10.2038 \\
16 & 1 & 4 & 34.81 & $1.6397\pm0.0436$ & 0.9925 & 0.0135 & 1.6852 & 3.4922 & 3.7368 & 5.5875 \\
16 & 1 & 8 & 69.52 & $1.5915\pm0.0058$ & 0.9925 & 0.0122 & 1.6161 & 4.0661 & 4.1889 & 8.8301 \\
16 & 2 & 4 & 34.81 & $1.5903\pm0.0045$ & 0.9927 & 0.0116 & 1.6448 & 5.5960 & 7.4676 & 12.5932 \\
16 & 2 & 8 & 69.52 & $1.5765\pm0.0172$ & 0.9919 & 0.0156 & 1.6064 & 5.9542 & 6.7282 & 10.2379 \\
\bottomrule
\end{tabular}}
\end{table*}

Top-2 improves dense loss in all six matched comparisons (mean paired change
$-1.75\%$), and eight layers improve all six depth comparisons (mean
$-5.05\%$). The best dense cell is 4 experts, top-2, and 8 layers at
$1.5606\pm0.0102$. These are fixed-budget comparisons, not scaling laws.

\begin{table}[t]
\centering
\scriptsize
\caption{Byte-tokenized dataset robustness for the 8-expert/top-1/4-layer
reference architecture. Dense uncertainty is sample SD; losses are not
comparable across corpora.}
\label{tab:app_dataset_grid}
\resizebox{\linewidth}{!}{%
\begin{tabular}{@{}lrrrrr@{}}
\toprule
Dataset & Dense & 50\% mag. & 50\% random & 80\% mag. & 80\% random \\
\midrule
TinyStories & $1.0206\pm0.0096$ & 1.0676 & 3.2700 & 3.5043 & 5.1919 \\
WikiText-103 & $1.6859\pm0.0160$ & 1.7388 & 4.0564 & 4.5586 & 9.3660 \\
Balanced mixture & $1.4801\pm0.0425$ & 1.5360 & 3.7565 & 5.5286 & 7.1896 \\
\bottomrule
\end{tabular}}
\end{table}

The completed representative Phase-4 rewind suite contains 288 retrainings:
three architecture/data representatives, three seeds, two sparsities, four
rewind points, and four conditions (a learned mask plus three controls). Each
retraining runs the full 2,500 updates.
\Cref{tab:app_phase4_rewind} reports the best learned-mask point per cell; its
``criterion'' column uses the manuscript's 5\% dense-loss tolerance.

\begin{table*}[t]
\centering
\scriptsize
\caption{Representative Phase-4 rewind summary. Recorded SDs are population
SDs. ``Controls'' means the learned mask beats random mask, random reinit, and
randomized routing at that rewind.}
\label{tab:app_phase4_rewind}
\resizebox{\textwidth}{!}{%
\begin{tabular}{@{}lrrrrcc@{}}
\toprule
Representative & Dense & Sparsity & Best rewind & Learned loss & Beats controls & $\leq5\%$ \\
\midrule
4E/top-2/8L WikiText & $1.5606\pm0.0083$ & 50\% & 1\% & $1.5654\pm0.0113$ & No & Yes \\
4E/top-2/8L WikiText & $1.5606\pm0.0083$ & 80\% & 10\% & $1.6385\pm0.0067$ & Yes & Yes \\
16E/top-1/8L WikiText & $1.5915\pm0.0047$ & 50\% & 0\% & $1.5846\pm0.0175$ & No & Yes \\
16E/top-1/8L WikiText & $1.5915\pm0.0047$ & 80\% & 10\% & $1.5710\pm0.0085$ & Yes & Yes \\
8E/top-1/4L mixture & $1.4801\pm0.0347$ & 50\% & 10\% & $1.4007\pm0.0253$ & Yes & Yes \\
8E/top-1/4L mixture & $1.4801\pm0.0347$ & 80\% & 10\% & $1.4792\pm0.0200$ & Yes & Yes \\
\bottomrule
\end{tabular}}
\end{table*}

The 50\% top-2 and high-capacity learned masks preserve dense loss but do not
beat every control at their selected points. All three 80\% representative
cells satisfy the 5\% loss criterion and beat the controls. Because archived
top-2 traces contain only primary IDs, these architecture results do not provide
a full top-2 replay analysis.

\section{Long-Budget Suite}
\label{app:long_budget}

The only completed 10,000-step family is the balanced-mixture, byte-tokenized,
8-expert/top-1/4-layer configuration with seeds 7, 17, and 29. Batch size 128
and sequence length 128 correspond to approximately $1.64\times10^8$ processed
tokens per run. The validation-batch cap is 64, but the original validation file
contains only about 31 batches and is exhausted. Saved checkpoints are
\begin{equation}
t\in\{0,100,500,1000,2500,5000,10000\}.
\label{eq:app_long_checkpoints}
\end{equation}
Validation is also logged every 500 steps. The final-step dense comparison is
misleading: learned routing improves through the middle of training and then
finishes worse on the same validation stream. \Cref{tab:app_long_checkpoint}
separates final, best observed, and best saved losses. ``Best'' is selected on
the same original validation stream and is therefore model-selection biased.

\begin{table}[t]
\centering
\scriptsize
\caption{Long-budget checkpoint comparison. Final uncertainty is sample SD;
best-observed and best-saved columns are three-seed means.}
\label{tab:app_long_checkpoint}
\resizebox{\linewidth}{!}{%
\begin{tabular}{@{}lrrrr@{}}
\toprule
Routing & Final 10k & Best observed & Best saved & $\Delta$ vs normal best saved \\
\midrule
Learned & $1.5236\pm0.0415$ & 1.3050 & 1.3072 & -- \\
Fixed random & $1.4175\pm0.0517$ & 1.3037 & 1.3289 & $+1.65\%$ \\
Random/step & $1.4313\pm0.0174$ & 1.4313 & 1.4313 & $+9.49\%$ \\
Shuffled usage & $1.4417\pm0.0086$ & 1.4417 & 1.4417 & $+10.29\%$ \\
\bottomrule
\end{tabular}}
\end{table}

The learned-routing best saved checkpoint is step 5,000 for every seed, with
per-seed losses 1.3046, 1.3376, and 1.2795. Fixed random also uses step 5,000;
random/step and shuffled usage use step 10,000. Direct pruning at these selected
checkpoints is reported in \cref{tab:app_long_pruning}.

\begin{table*}[t]
\centering
\scriptsize
\caption{Long-budget direct pruning at each condition's best saved checkpoint.
Entries are three-seed means; uncertainty is omitted for table width.}
\label{tab:app_long_pruning}
\resizebox{\textwidth}{!}{%
\begin{tabular}{@{}lrrrrrrrrr@{}}
\toprule
Routing & Dense & 50\% mag. & 50\% random & 50\% transfer & 50\% reinit & 80\% mag. & 80\% random & 80\% transfer & 80\% reinit \\
\midrule
Learned & 1.3072 & 1.3905 & 5.7336 & 1.8534 & 19.8493 & 7.7513 & 15.5672 & 9.0382 & 19.7335 \\
Fixed random & 1.3289 & 1.4003 & 5.3759 & 1.7121 & 24.3530 & 7.8883 & 19.9285 & 8.4725 & 24.6300 \\
Random/step & 1.4313 & 1.4835 & 3.0754 & 1.6340 & 5.6137 & 2.6843 & 5.0693 & 2.9500 & 5.6110 \\
Shuffled usage & 1.4417 & 1.4593 & 1.9987 & 1.5046 & 2.8512 & 1.7413 & 2.6652 & 1.8273 & 2.8424 \\
\bottomrule
\end{tabular}}
\end{table*}

At the final 10,000-step checkpoints, the per-seed Pearson correlation between
final validation-route agreement and 80\% mask Jaccard over the six unordered
pairs among four routing conditions is $0.4288\pm0.0486$. This descriptive
association is weaker than in the original single-domain suites and is not used
for checkpoint selection.

The reduced rewind follow-up extracts only 80\% masks from each seed's best
saved learned-routing checkpoint. It evaluates the learned mask at 0\%, 10\%,
and 25\% rewind; matched controls were run only at the winning initialization
point. Every lane retrains for a full 10,000 updates. Contrary to the
2,500-step suites, initialization is best here.

\begin{table}[t]
\centering
\scriptsize
\caption{Reduced long-budget 80\% rewind suite, mean $\pm$ population SD. The
best-saved dense reference is $1.3072\pm0.0238$.}
\label{tab:app_long_rewind}
\begin{tabular}{@{}lrrr@{}}
\toprule
Condition & Rewind & Loss & $\Delta$ vs dense \\
\midrule
Learned mask & 0\% & $1.2984\pm0.0055$ & $-0.68\%$ \\
Learned mask & 10\% & $1.3552\pm0.0087$ & $+3.67\%$ \\
Learned mask & 25\% & $1.4125\pm0.0650$ & $+8.05\%$ \\
Random mask & 0\% & $1.3110\pm0.0070$ & $+0.29\%$ \\
Random reinit & 0\% & $1.3088\pm0.0074$ & $+0.12\%$ \\
Randomized routing & 0\% & $1.5011\pm0.0129$ & $+14.83\%$ \\
\bottomrule
\end{tabular}
\end{table}

The learned mask's margins over the random mask and random reinitialization are
small; the randomized-routing gap is much larger. An expanded validation pass
uses 8,740 sequences (about 69 batches) assembled from the existing TinyStories
validation source plus WikiText validation and test text. It overlaps the
original validation sources and is not an independent selection split, but it
increases coverage. On this pass, dense, learned-mask, random-mask, and
random-reinit losses are respectively $1.3213\pm0.0203$,
$1.3096\pm0.0044$, $1.3242\pm0.0080$, and $1.3225\pm0.0099$. The narrow
learned-mask advantage remains, with the same caveat about checkpoint and mask
selection. Uncertainty in this expanded-validation comparison is population SD.

\section{Additional Completed Controls}
\label{app:additional_controls}

\subsection{Load-Balancing Sweep}

The load-balancing sweep varies only
$\lambda_{\mathrm{bal}}\in\{0,0.01,0.03,0.1,0.3\}$ within each dataset and
seed. \Cref{tab:app_load_balance} shows that the principal 0.1 setting avoids
the collapse seen at zero, while 0.3 gives the best dense mean. Magnitude masks
beat random masks at every tested coefficient, although one-shot 80\% pruning
becomes more brittle as balancing strengthens.

An earlier single-seed, two-layer CPU pilot is retained only as motivation for
this sweep. At coefficient 0.01 its two layers ended with normalized usage
entropy 0.506 and 0.194 and one dead expert each; at coefficient 0.1 entropy
rose to 0.968 and 0.979 with no dead experts. Because its architecture and
replication differ from the GPU suites, these pilot values are not pooled with
the principal results.

\begin{table*}[t]
\centering
\scriptsize
\caption{Load-balance sweep, mean over three seeds. Dense uncertainty is sample
SD. Selected-route entropy $H$ and dead experts are final logged-step values
averaged over layers within each run and then seeds; direct-pruning losses are
also recorded.}
\label{tab:app_load_balance}
\resizebox{\textwidth}{!}{%
\begin{tabular}{@{}llrrrrrrr@{}}
\toprule
Dataset & $\lambda$ & Dense & $H$ & Dead & 50\% mag. & 50\% random & 80\% mag. & 80\% random \\
\midrule
TinyStories & 0 & $1.2427\pm0.0142$ & 0.5112 & 1.08 & 1.2747 & 1.9416 & 1.7854 & 2.4443 \\
 & 0.01 & $1.0634\pm0.0117$ & 0.9610 & 0 & 1.1073 & 2.4795 & 2.2598 & 3.6078 \\
 & 0.03 & $1.0439\pm0.0166$ & 0.9856 & 0 & 1.0882 & 2.7105 & 2.5058 & 4.0400 \\
 & 0.1 & $1.0206\pm0.0096$ & 0.9957 & 0 & 1.0676 & 3.2700 & 3.5043 & 5.1919 \\
 & 0.3 & $1.0146\pm0.0145$ & 0.9971 & 0 & 1.0647 & 4.2315 & 5.7921 & 8.8968 \\
\midrule
WikiText-103 & 0 & $1.9252\pm0.0087$ & 0.2535 & 3.83 & 1.9439 & 2.3248 & 2.2530 & 2.6468 \\
 & 0.01 & $1.7231\pm0.0249$ & 0.9582 & 0 & 1.7622 & 2.7096 & 2.7351 & 3.4988 \\
 & 0.03 & $1.7040\pm0.0280$ & 0.9760 & 0 & 1.7474 & 3.0594 & 3.8592 & 5.1037 \\
 & 0.1 & $1.6859\pm0.0160$ & 0.9932 & 0 & 1.7388 & 4.0564 & 4.5586 & 9.3660 \\
 & 0.3 & $1.6502\pm0.0234$ & 0.9970 & 0 & 1.7128 & 5.4444 & 7.1457 & 15.6605 \\
\bottomrule
\end{tabular}}
\end{table*}

\subsection{Routing-Condition Pruning and Fixed-Router Rewind}

\Cref{tab:app_multidomain_pruning} supplies the full numeric pruning comparison
behind the byte and byte-ngram causal-control discussion. Coherent learned and
fixed-random routing yield a large magnitude--random-mask gap; random/step and
shuffled routing are much less sensitive to mask identity.

\begin{table*}[t]
\centering
\scriptsize
\caption{Balanced-mixture direct pruning by tokenizer and routing condition,
mean loss over three seeds. Sample SD is omitted for table width.}
\label{tab:app_multidomain_pruning}
\resizebox{\textwidth}{!}{%
\begin{tabular}{@{}llrrrrrrrrr@{}}
\toprule
Tokenizer & Routing & Dense & 50\% mag. & 50\% random & 50\% transfer & 50\% reinit & 80\% mag. & 80\% random & 80\% transfer & 80\% reinit \\
\midrule
Byte & Learned & 1.4801 & 1.5360 & 3.7565 & 1.8653 & 8.2341 & 5.5286 & 7.1896 & 5.9469 & 8.3213 \\
 & Fixed random & 1.5810 & 1.6244 & 3.1311 & 1.8658 & 8.3148 & 4.8987 & 6.4320 & 5.5288 & 8.3894 \\
 & Random/step & 1.9564 & 1.9788 & 2.3928 & 2.0219 & 2.8853 & 2.3164 & 2.7988 & 2.3687 & 2.8852 \\
 & Shuffled usage & 1.9791 & 1.9819 & 2.0373 & 1.9877 & 2.0841 & 2.0142 & 2.0775 & 2.0212 & 2.0841 \\
\midrule
Byte-ngram & Learned & 3.2753 & 3.3992 & 11.9938 & 3.7630 & 39.4429 & 14.0545 & 35.6571 & 17.4116 & 39.6119 \\
 & Fixed random & 3.5037 & 3.6076 & 14.8263 & 4.1261 & 34.2940 & 17.3239 & 31.3733 & 21.4662 & 34.7452 \\
 & Random/step & 4.3504 & 4.3601 & 4.5309 & 4.3778 & 4.8031 & 4.4818 & 4.7644 & 4.5099 & 4.8030 \\
 & Shuffled usage & 4.3574 & 4.3623 & 4.4146 & 4.3675 & 4.4898 & 4.4112 & 4.4794 & 4.4161 & 4.4898 \\
\bottomrule
\end{tabular}}
\end{table*}

Fixed-random routing freezes only the router projection; its assignments can
still change as hidden states evolve. On WikiText its dense loss is
$1.7386\pm0.0206$ (population SD in this follow-up). Both learned-mask curves
are best at the full-budget 10\% rewind; \cref{tab:app_fixed_router_rewind}
reports all matched controls there.

\begin{table}[t]
\centering
\scriptsize
\caption{Fixed-random-router full-budget retraining at the 10\% rewind, mean
$\pm$ population SD over three seeds.}
\label{tab:app_fixed_router_rewind}
\resizebox{\linewidth}{!}{%
\begin{tabular}{@{}rrrrrr@{}}
\toprule
Sparsity & Dense & Learned mask & Random mask & Random reinit & Randomized routing \\
\midrule
50\% & $1.7386\pm0.0206$ & $1.6315\pm0.0167$ & $1.7334\pm0.0206$ & $1.7390\pm0.0394$ & $2.0863\pm0.0158$ \\
80\% & $1.7386\pm0.0206$ & $1.7165\pm0.0042$ & $1.8899\pm0.0452$ & $1.8653\pm0.0317$ & $2.1022\pm0.0081$ \\
\bottomrule
\end{tabular}}
\end{table}

The learned-mask gaps are $-6.16\%$ and $-1.27\%$ at 50\% and 80\%,
respectively. Thus learned router adaptation is useful
for the best dense performance but is not required for a rewindable sparse
support under this protocol.

\subsection{Temporal Stability and Expert Diagnostics}

The repository also contains representative seed-7 checkpoint analyses that
are not treated as independent replications. Step-25-to-step-2,500 same-token
primary-route agreement is low overall but above the $1/8$ uniform baseline,
with the first layer most stable (\cref{tab:app_temporal_routes}).

\begin{table}[t]
\centering
\scriptsize
\caption{Representative seed-7 learned-routing step-25-to-step-2,500
validation-route agreement.}
\label{tab:app_temporal_routes}
\begin{tabular}{@{}lrrrrr@{}}
\toprule
Dataset & Layer 0 & Layer 1 & Layer 2 & Layer 3 & Overall \\
\midrule
TinyStories & 0.3112 & 0.1226 & 0.1383 & 0.1380 & 0.1776 \\
WikiText-103 & 0.3394 & 0.0890 & 0.1604 & 0.1420 & 0.1827 \\
\bottomrule
\end{tabular}
\end{table}

The WikiText checkpoint analysis additionally compares each checkpoint's
80\%-sparse expert-local magnitude mask with its own final support. Jaccard
increases gradually under both routing regimes (\cref{tab:app_temporal_masks});
the support is not fully fixed at initialization. Random-routing masks become
more similar to their final mask earlier, but those runs have worse dense and
rewind performance, so temporal stability alone is not evidence of a useful
ticket.

\begin{table*}[t]
\centering
\scriptsize
\caption{Representative seed-7 WikiText 80\% mask Jaccard with the final mask.}
\label{tab:app_temporal_masks}
\begin{tabular}{@{}lrrrrrrr@{}}
\toprule
Routing & Step 0 & 25 & 125 & 250 & 500 & 1,000 & 2,500 \\
\midrule
Learned & 0.5795 & 0.5969 & 0.6023 & 0.6076 & 0.6199 & 0.6518 & 1.0000 \\
Random/step & 0.5147 & 0.5637 & 0.6840 & 0.6964 & 0.7267 & 0.7883 & 1.0000 \\
\bottomrule
\end{tabular}
\end{table*}

Final-checkpoint expert diagnostics are summarized in
\cref{tab:app_expert_diagnostics}. A linear probe predicts the selected primary
expert from pre-router hidden states, while the silhouette score measures
geometric cluster separation. The probe uses at most 5,000 positions, a
stratified 75/25 train/test split, and logistic regression with
\texttt{max\_iter=200} and \texttt{random\_state=0}. The substitution analysis
reroutes the token positions selected for each source expert to each candidate
expert one layer at a time and compares token cross entropy; the reported
column counts source-expert subsets for which their own expert has the lowest
substituted loss.
Prediction is strong, especially in early layers, but low or negative
silhouette scores and weaker late-layer substitution diagonals caution against
equating route predictability with clean clusters or semantic specialization.
Moreover, the router itself selects the argmax of a linear projection of these
same hidden states, so high linear-probe accuracy is structurally expected and
is not independent evidence of specialization. The WikiText probe emitted
logistic-regression convergence warnings at the recorded iteration limit, so
its accuracies are approximate.

\begin{table}[t]
\centering
\scriptsize
\caption{Representative seed-7 final learned-routing checkpoint diagnostics.
``Own best'' is out of eight experts.}
\label{tab:app_expert_diagnostics}
\begin{tabular}{@{}lrrrr@{}}
\toprule
Dataset & Layer & Probe accuracy & Silhouette & Own best \\
\midrule
TinyStories & 0 & 0.9936 & 0.0514 & 8/8 \\
TinyStories & 1 & 0.9544 & 0.0451 & 8/8 \\
TinyStories & 2 & 0.8664 & $-0.0149$ & 7/8 \\
TinyStories & 3 & 0.7952 & $-0.0198$ & 3/8 \\
WikiText-103 & 0 & 0.9960 & 0.0563 & 8/8 \\
WikiText-103 & 1 & 0.9176 & 0.0103 & 8/8 \\
WikiText-103 & 2 & 0.8208 & $-0.0516$ & 4/8 \\
WikiText-103 & 3 & 0.7936 & $-0.0235$ & 2/8 \\
\bottomrule
\end{tabular}
\end{table}

\subsection{Representative Iterative Magnitude Pruning}

The completed IMP check uses the balanced-mixture 8E/top-1/4L cell, seeds 7,
17, and 29, and initialization rewind. Each of four rounds prunes currently
surviving expert weights, rewinds the survivors, and retrains for a full 2,500
updates. Results are in \cref{tab:app_imp}.

\begin{table}[t]
\centering
\scriptsize
\caption{Representative initialization-rewind IMP, mean $\pm$ population SD.
The dense reference is $1.4801\pm0.0347$.}
\label{tab:app_imp}
\begin{tabular}{@{}rrrr@{}}
\toprule
Round & Sparsity & Loss & $\Delta$ vs dense \\
\midrule
1 & 33.13\% & $1.4938\pm0.0060$ & $+0.93\%$ \\
2 & 55.28\% & $1.5013\pm0.0312$ & $+1.44\%$ \\
3 & 70.09\% & $1.5168\pm0.0445$ & $+2.48\%$ \\
4 & 80.00\% & $1.5687\pm0.0365$ & $+5.99\%$ \\
\bottomrule
\end{tabular}
\end{table}

At 80\%, IMP improves over the one-shot initialization-restart result
$1.6549\pm0.0177$ in this cell, but it is worse than the one-shot mask with a
10\% rewind, $1.4792\pm0.0200$. This negative comparison supports using the
simpler one-shot-plus-early-rewind protocol for the principal high-sparsity
claim; it does not establish that IMP is generally inferior.

\end{document}
